\section{Cronograma de actividades}
\label{sec:cronograma}

El experimento de enseñanza se organiza en tres fases principales: diseño, implementación y análisis. La implementación en aula abarca 10 semanas; el análisis de datos y la escritura de resultados se desarrollan en las semanas siguientes. El cronograma a continuación refleja la planificación de actividades para el primer año de la maestría.

\begin{table}[htbp]
\centering
\caption{Cronograma de actividades del experimento de enseñanza}
\label{tab:cronograma}
\small
\begin{tabular}{|p{5.5cm}|c|c|c|c|c|c|c|c|c|c|}
\hline
\textbf{Actividad} & \textbf{S1} & \textbf{S2} & \textbf{S3} & \textbf{S4} & \textbf{S5} & \textbf{S6} & \textbf{S7} & \textbf{S8} & \textbf{S9} & \textbf{S10} \\
\hline
\multicolumn{11}{|l|}{\textit{Fase 1: Diseño}} \\
\hline
Revisión de literatura y ajuste del marco teórico & $\bullet$ & $\bullet$ & & & & & & & & \\
\hline
Diseño definitivo de tareas 1 y 2 & $\bullet$ & $\bullet$ & & & & & & & & \\
\hline
Construcción de instrumentos (guía obs., protocolo entrevistas, rúbrica) & $\bullet$ & $\bullet$ & & & & & & & & \\
\hline
Trámite de consentimientos informados & $\bullet$ & $\bullet$ & & & & & & & & \\
\hline
Configuración y prueba de herramientas (GeoGebra + ChatGPT) & & $\bullet$ & $\bullet$ & & & & & & & \\
\hline
\multicolumn{11}{|l|}{\textit{Fase 2: Implementación}} \\
\hline
Pre-test y sesión introductoria & & & $\bullet$ & & & & & & & \\
\hline
Tarea 1 — La rueda de la fortuna (2 sesiones, 40 min c/u) & & & & $\bullet$ & $\bullet$ & & & & & \\
\hline
Primeras entrevistas con muestra intensiva & & & & & $\bullet$ & & & & & \\
\hline
Tarea 2 — El movimiento de las mareas (2 sesiones, 40 min c/u) & & & & & & $\bullet$ & $\bullet$ & & & \\
\hline
Segundas entrevistas con muestra intensiva & & & & & & & $\bullet$ & & & \\
\hline
Post-test y entrevistas finales & & & & & & & & $\bullet$ & & \\
\hline
\multicolumn{11}{|l|}{\textit{Fase 3: Análisis}} \\
\hline
Organización del corpus: transcripciones, digitalización & & & & & & & & $\bullet$ & $\bullet$ & \\
\hline
Codificación abierta y axial & & & & & & & & & $\bullet$ & $\bullet$ \\
\hline
Análisis de patrones y triangulación & & & & & & & & & & $\bullet$ \\
\hline
\end{tabular}
\end{table}

\vspace{0.5em}

Las semanas 1--2 corresponden a la fase de preparación previa al inicio de la implementación en aula. Las semanas 3--8 constituyen el período de implementación activa. Las semanas 8--10 inician el análisis, que se extiende en el segundo semestre del año hasta la producción de los resultados y la escritura del informe final.

Este cronograma es indicativo. El experimento de enseñanza, por su naturaleza iterativa, admite ajustes en la secuencia cuando los datos recolectados lo justifiquen.
